\documentclass[../main.tex]{subfiles}

\begin{document}
\begin{problema}
    Dada \(u \colon \Omega \subseteq \mathbb{C} \to \mathbb{R}\), definimos
    \(\tilde{u} \colon \tilde{\Omega} \to \mathbb{R}\) como \(\tilde{u}(z) =
    u(\overline{z})\), donde \(\tilde{\Omega}\) está definido como en el ejercicio
    anterior. Pruebe que \(u\) es armónica en \(\Omega\) si y solo si \(\tilde{u}\) es
    armónica en \(\tilde{\Omega}\).
\end{problema}
\end{document}
